% Options for packages loaded elsewhere
\PassOptionsToPackage{unicode}{hyperref}
\PassOptionsToPackage{hyphens}{url}
%
\documentclass[
  12pt,
  oneside]{article}
\author{}
\date{\vspace{-2.5em}}

\usepackage{amsmath,amssymb}
\usepackage{lmodern}
\usepackage{setspace}
\usepackage{iftex}
\ifPDFTeX
  \usepackage[T1]{fontenc}
  \usepackage[utf8]{inputenc}
  \usepackage{textcomp} % provide euro and other symbols
\else % if luatex or xetex
  \usepackage{unicode-math}
  \defaultfontfeatures{Scale=MatchLowercase}
  \defaultfontfeatures[\rmfamily]{Ligatures=TeX,Scale=1}
\fi
% Use upquote if available, for straight quotes in verbatim environments
\IfFileExists{upquote.sty}{\usepackage{upquote}}{}
\IfFileExists{microtype.sty}{% use microtype if available
  \usepackage[]{microtype}
  \UseMicrotypeSet[protrusion]{basicmath} % disable protrusion for tt fonts
}{}
\makeatletter
\@ifundefined{KOMAClassName}{% if non-KOMA class
  \IfFileExists{parskip.sty}{%
    \usepackage{parskip}
  }{% else
    \setlength{\parindent}{0pt}
    \setlength{\parskip}{6pt plus 2pt minus 1pt}}
}{% if KOMA class
  \KOMAoptions{parskip=half}}
\makeatother
\usepackage{xcolor}
\IfFileExists{xurl.sty}{\usepackage{xurl}}{} % add URL line breaks if available
\IfFileExists{bookmark.sty}{\usepackage{bookmark}}{\usepackage{hyperref}}
\hypersetup{
  hidelinks,
  pdfcreator={LaTeX via pandoc}}
\urlstyle{same} % disable monospaced font for URLs
\usepackage[top=1in,bottom=1in,right=1.1in,left=1.1in]{geometry}
\usepackage{color}
\usepackage{fancyvrb}
\newcommand{\VerbBar}{|}
\newcommand{\VERB}{\Verb[commandchars=\\\{\}]}
\DefineVerbatimEnvironment{Highlighting}{Verbatim}{commandchars=\\\{\}}
% Add ',fontsize=\small' for more characters per line
\newenvironment{Shaded}{}{}
\newcommand{\AlertTok}[1]{\textcolor[rgb]{1.00,0.00,0.00}{\textbf{#1}}}
\newcommand{\AnnotationTok}[1]{\textcolor[rgb]{0.38,0.63,0.69}{\textbf{\textit{#1}}}}
\newcommand{\AttributeTok}[1]{\textcolor[rgb]{0.49,0.56,0.16}{#1}}
\newcommand{\BaseNTok}[1]{\textcolor[rgb]{0.25,0.63,0.44}{#1}}
\newcommand{\BuiltInTok}[1]{#1}
\newcommand{\CharTok}[1]{\textcolor[rgb]{0.25,0.44,0.63}{#1}}
\newcommand{\CommentTok}[1]{\textcolor[rgb]{0.38,0.63,0.69}{\textit{#1}}}
\newcommand{\CommentVarTok}[1]{\textcolor[rgb]{0.38,0.63,0.69}{\textbf{\textit{#1}}}}
\newcommand{\ConstantTok}[1]{\textcolor[rgb]{0.53,0.00,0.00}{#1}}
\newcommand{\ControlFlowTok}[1]{\textcolor[rgb]{0.00,0.44,0.13}{\textbf{#1}}}
\newcommand{\DataTypeTok}[1]{\textcolor[rgb]{0.56,0.13,0.00}{#1}}
\newcommand{\DecValTok}[1]{\textcolor[rgb]{0.25,0.63,0.44}{#1}}
\newcommand{\DocumentationTok}[1]{\textcolor[rgb]{0.73,0.13,0.13}{\textit{#1}}}
\newcommand{\ErrorTok}[1]{\textcolor[rgb]{1.00,0.00,0.00}{\textbf{#1}}}
\newcommand{\ExtensionTok}[1]{#1}
\newcommand{\FloatTok}[1]{\textcolor[rgb]{0.25,0.63,0.44}{#1}}
\newcommand{\FunctionTok}[1]{\textcolor[rgb]{0.02,0.16,0.49}{#1}}
\newcommand{\ImportTok}[1]{#1}
\newcommand{\InformationTok}[1]{\textcolor[rgb]{0.38,0.63,0.69}{\textbf{\textit{#1}}}}
\newcommand{\KeywordTok}[1]{\textcolor[rgb]{0.00,0.44,0.13}{\textbf{#1}}}
\newcommand{\NormalTok}[1]{#1}
\newcommand{\OperatorTok}[1]{\textcolor[rgb]{0.40,0.40,0.40}{#1}}
\newcommand{\OtherTok}[1]{\textcolor[rgb]{0.00,0.44,0.13}{#1}}
\newcommand{\PreprocessorTok}[1]{\textcolor[rgb]{0.74,0.48,0.00}{#1}}
\newcommand{\RegionMarkerTok}[1]{#1}
\newcommand{\SpecialCharTok}[1]{\textcolor[rgb]{0.25,0.44,0.63}{#1}}
\newcommand{\SpecialStringTok}[1]{\textcolor[rgb]{0.73,0.40,0.53}{#1}}
\newcommand{\StringTok}[1]{\textcolor[rgb]{0.25,0.44,0.63}{#1}}
\newcommand{\VariableTok}[1]{\textcolor[rgb]{0.10,0.09,0.49}{#1}}
\newcommand{\VerbatimStringTok}[1]{\textcolor[rgb]{0.25,0.44,0.63}{#1}}
\newcommand{\WarningTok}[1]{\textcolor[rgb]{0.38,0.63,0.69}{\textbf{\textit{#1}}}}
\usepackage{longtable,booktabs,array}
\usepackage{calc} % for calculating minipage widths
% Correct order of tables after \paragraph or \subparagraph
\usepackage{etoolbox}
\makeatletter
\patchcmd\longtable{\par}{\if@noskipsec\mbox{}\fi\par}{}{}
\makeatother
% Allow footnotes in longtable head/foot
\IfFileExists{footnotehyper.sty}{\usepackage{footnotehyper}}{\usepackage{footnote}}
\makesavenoteenv{longtable}
\usepackage{graphicx}
\makeatletter
\def\maxwidth{\ifdim\Gin@nat@width>\linewidth\linewidth\else\Gin@nat@width\fi}
\def\maxheight{\ifdim\Gin@nat@height>\textheight\textheight\else\Gin@nat@height\fi}
\makeatother
% Scale images if necessary, so that they will not overflow the page
% margins by default, and it is still possible to overwrite the defaults
% using explicit options in \includegraphics[width, height, ...]{}
\setkeys{Gin}{width=\maxwidth,height=\maxheight,keepaspectratio}
% Set default figure placement to htbp
\makeatletter
\def\fps@figure{htbp}
\makeatother
\setlength{\emergencystretch}{3em} % prevent overfull lines
\providecommand{\tightlist}{%
  \setlength{\itemsep}{0pt}\setlength{\parskip}{0pt}}
\setcounter{secnumdepth}{-\maxdimen} % remove section numbering
\usepackage{booktabs}
\usepackage{longtable}
\usepackage{array}
\usepackage{multirow}
\usepackage{wrapfig}
\usepackage{float}
\usepackage{colortbl}
\usepackage{pdflscape}
\usepackage{tabu}
\usepackage{threeparttable}
\usepackage{threeparttablex}
\usepackage[normalem]{ulem}
\usepackage{makecell}
\usepackage{xcolor}
\ifLuaTeX
  \usepackage{selnolig}  % disable illegal ligatures
\fi

\begin{document}

%%%%%%%%%%%%%%%%%%%%%%%%%%%%%%%%%%%%%%%%%
% University Assignment Title Page
% LaTeX Template
% Version 1.0 (27/12/12)
%
% This template has been downloaded from:
% http://www.LaTeXTemplates.com
%
% Original author:
% WikiBooks (http://en.wikibooks.org/wiki/LaTeX/Title_Creation)
%
% License:
% CC BY-NC-SA 3.0 (http://creativecommons.org/licenses/by-nc-sa/3.0/)

%-------------------------------------------------------------------------------
%	PACKAGES AND OTHER DOCUMENT CONFIGURATIONS
%-------------------------------------------------------------------------------


\begin{titlepage}

\newcommand{\HRule}{\rule{\linewidth}{0.5mm}} % Defines a new command for the horizontal lines, change thickness here

\center % Center everything on the page

%----------------------------------------------------------------------------------------
%	HEADING SECTIONS
%----------------------------------------------------------------------------------------
%
% \textsc{\LARGE Institute of Statistical Reasearch}\\[0.3cm] % Name of your university/college
% \textsc{\LARGE and Training}\\[0.3cm]
% \textsc{\LARGE University of Dhaka}\\[0.6cm]
% Minor heading such as course title

%-------------------------------------------------------------------------------
%	TITLE SECTION
%-------------------------------------------------------------------------------

\begin{minipage}{3cm}
\vspace{3cm}
\end{minipage}


\HRule \\[0.25cm]
{ \huge \bfseries Assignment}\\[0.06cm]
{ \huge \bfseries On}\\[0.1cm]
{ \huge \bfseries Lifetime Data Analysis}\\[0.08cm]
{ \LARGE \bfseries AST 405}\\[0.05cm]
\HRule \\[1.5cm]

%-------------------------------------------------------------------------------
%	AUTHOR SECTION
%-------------------------------------------------------------------------------

\begin{minipage}{0.4\textwidth}
\begin{flushleft} \large
\emph{Submitted By:}\\
Shafayet Khan Shafee \\FH-033-011\\4th Year
\end{flushleft}
\end{minipage}
~
\begin{minipage}{0.4\textwidth}
\begin{flushright} \large
\emph{Submitted To:} \\
Dr. Mahbub Latif\\Professor\\ISRT, DU
\end{flushright}
\end{minipage}\\[1cm]

%-------------------------------------------------------------------------------
%	DATE SECTION
%-------------------------------------------------------------------------------

{\large \today}\\[3cm]

%-------------------------------------------------------------------------------
%	LOGO SECTION
%-------------------------------------------------------------------------------

\begin{figure}[hb]
  \centering
  \includegraphics[width=4cm, height=5.2cm]{../others/du_logo_high.JPG}
\end{figure}


%-------------------------------------------------------------------------------

\vfill

\end{titlepage}

{
\setcounter{tocdepth}{2}
\tableofcontents
}
\setstretch{1.25}
\clearpage

\hypertarget{answer-to-ques-01}{%
\subsection{Answer To Ques 01}\label{answer-to-ques-01}}

For the continuous variables \texttt{age} and \texttt{bmi}, mean and standard deviation (SD) is obtained. And for categorical variable \texttt{arcus}, \texttt{behpat} and \texttt{chd69}, frequency and proportion (in percentage) is obtained as descriptive statistics, which are shown in the following table:

\begin{table}[!h]

\caption{\label{tab:Table-01}Descriptive Statistics}
\centering
\begin{tabular}[t]{lll}
\toprule
  & level & Overall\\
\midrule
n &  & 3154\\
age (mean (SD)) &  & 46.35 (5.56)\\
bmi (mean (SD)) &  & 24.48 (2.55)\\
arcus (\%) & 0 & 2219 (70.4)\\
 & 1 & 934 (29.6)\\
\addlinespace
behpat (\%) & A1 & 275 ( 8.7)\\
 & A2 & 1290 (40.9)\\
 & B3 & 1236 (39.2)\\
 & B4 & 353 (11.2)\\
chd69 (\%) & No & 2900 (91.9)\\
\addlinespace
 & Yes & 254 ( 8.1)\\
\bottomrule
\end{tabular}
\end{table}

\clearpage

\hypertarget{answer-to-ques-02}{%
\subsection{Answer To Ques 02}\label{answer-to-ques-02}}

\hypertarget{effect-of-age-and-bmi}{%
\paragraph{Effect of age and bmi :}\label{effect-of-age-and-bmi}}

To check whether each of \texttt{age} and \texttt{bmi} has significant effect on \texttt{chol}, we need to do correlation test (pearson).

For each correlation test of \texttt{age} and \texttt{bmi} with \texttt{chol}, the hypotheses are,
\begin{align*}
H_o&: \rho = 0 \\
H_a&: \rho \neq 0
\end{align*}
where \(\rho\) is the population correlation coefficient.

The appropriate test statistic for testing the hypothesis is,

\[
t_o = \frac{r\sqrt{n-2}}{\sqrt{1-r^2}}
\]

where \(r\) is the sample correlation coefficient and \(t_o\) follows \(t_{n-2}\)
distribution under \(H_o\).

Here we will reject \(H_o\) if the associated \texttt{p.value} is less then \(\alpha = 0.05\), where \(\alpha\) is the level of significance.

\clearpage

Now from Table: \ref{tab:Table-02}, we can see that both of the variables \texttt{age} and \texttt{bmi} is significantly associated with \texttt{chol} variable.

\begin{table}[!h]

\caption{\label{tab:Table-02}Correlation Test (age, bmi) and T-test (arcus) with chol}
\centering
\begin{tabular}[t]{lrrr}
\toprule
term & estimate & statistic & p.value\\
\midrule
age & 0.081 & 4.550 & <.001\\
bmi & 0.060 & 3.385 & <.001\\
arcus & -12.672 & -7.604 & <.001\\
\bottomrule
\end{tabular}
\end{table}

\hypertarget{effect-of-arcus}{%
\paragraph{Effect of arcus}\label{effect-of-arcus}}

Now to check the effect of \texttt{arcus} (a categorical variable with level 0 and 1) on \texttt{chol}, we can do two sample t-test, where one sample is people with \texttt{arcus} 0 and
other sample is people with \texttt{arcus} 1. Therefore, is \(\mu_o\) and \(\mu_1\) are the
population mean of \texttt{arcus} 0 and \texttt{arcus} 1 group respectively.

Then our hypotheses are:
\begin{align*}
H_o&: \mu_o = \mu_1 \\
H_a&: \mu_o \neq \mu_1
\end{align*}
and appropriate test statistic (assuming unequal variance),

\[
t_o = \frac{\bar{x_o} - \bar{x_1}}{\sqrt{\frac{s_o^2}{n_o} + \frac{s_1^2}{n_1}}}
\]

where \(\bar{x}\) and \(s^2\) denotes sample mean and sample variance with subscript
0 and 1 for \texttt{arcus} group 0 and 1 respectively and \(n_o\) and \(n_1\) are the corresponding sample sizes.

Here we will reject \(H_o\) if the associated \texttt{p\ value} is less then \(\alpha = 0.05\), where \(\alpha\) is the level of significance.

\clearpage

Then from Table: \ref{tab:Table-02}, since \texttt{p.value} is less than 0.001 we can conclude that mean \texttt{chol} differs significantly for 0 and 1 group of \texttt{arcus}, that is, \texttt{arcus} has significant effect on \texttt{chol}.

\hypertarget{effect-of-behpat}{%
\paragraph{Effect of behpat:}\label{effect-of-behpat}}

To check the effect of \texttt{behpat} on \texttt{chol}, we can do oneway ANOVA. In this case the hypotheses are:
\begin{align*}
H_o&: \mu_1 = \mu_2 \\
H_a&: \mu_i \neq \mu_j \! &\text{for at least on i $\ne$ j}
\end{align*}
and test statistic is \(F_o = \frac{MS_{reg}}{MS_{E}}\). We will reject \(H_o\) if associated \texttt{p.value} is less than 0.05.

\begin{table}[H]

\caption{\label{tab:Table-03}One Way Analysis of Variance for chol on behpat}
\centering
\begin{tabular}[t]{lrrrrr}
\toprule
term & df & sumsq & meansq & statistic & p.value\\
\midrule
behpat & 3 & 30741.67 & 10247.223 & 5.475 & <.001\\
Residuals & 3135 & 5867632.50 & 1871.653 & NA & NA\\
\bottomrule
\end{tabular}
\end{table}

Then from Table: \ref{tab:Table-03}, we conclude that mean \texttt{chol} differs significantly over the levels of \texttt{behpat}.

\clearpage

\hypertarget{answer-to-ques-03}{%
\subsection{Answer To Ques 03}\label{answer-to-ques-03}}

To examine the association between \texttt{behat} and \texttt{chd69}, the hypotheses are:
\begin{align*}
H_o&: \text{There's no association between \texttt{behpat} and \texttt{chd69}.} \\
H_a&: \text{There's association between \texttt{behpat} and \texttt{chd69}.}
\end{align*}
Here, the test statistics is:
\[
\chi^2 = \sum_{i=1}^{r}\sum_{j=1}^{c}\frac{(O_{ij} - E_{ij})^2}{E_{ij}}
\]
Where, \(O_{ij}\) is observed frequency and \(E_{ij}\) is the expected frequency and \(r\), \(c\) is the row and column number of contingency table \ref{Table-04}

Here, we would reject \(H_o\) if associated \texttt{p.value} corresponding to test statistics
is less the 0.05.

\begin{table}[H]
      \begin{minipage}{.5\linewidth}
      \centering
\begin{tabular}{lrr}
\toprule
\multicolumn{1}{c}{ } & \multicolumn{2}{c}{CHD} \\
\cmidrule(l{3pt}r{3pt}){2-3}
behpat & No & Yes\\
\midrule
A1 & 246 & 29\\
A2 & 1158 & 132\\
B3 & 1164 & 72\\
B4 & 332 & 21\\
\bottomrule
\end{tabular}\caption{Contingency table}
      \label{Table-04}
    \end{minipage}%
      \begin{minipage}{.5\linewidth}
        \centering
\begin{tabular}{rl}
\toprule
statistic & p.value\\
\midrule
20.978 & <.001\\
\bottomrule
\end{tabular}\caption{chi-square test}
    \label{Table-05}
    \end{minipage}
      \end{table}

Now from the Table \ref{Table-05}, since \texttt{p.value} is less than 0.001, we can conclude that there's an significant association between \texttt{behpat} and \texttt{chd69}.

\clearpage

\hypertarget{answer-to-ques-04}{%
\subsection{Answer To Ques 04}\label{answer-to-ques-04}}

\begin{table}[H]

\caption{\label{tab:Table-06}Estimate of Model Parameters}
\centering
\begin{tabular}[t]{lrrrr}
\toprule
term & estimate & std.error & statistic & p.value\\
\midrule
(Intercept) & 196.971 & 6.483 & 30.381 & <.001\\
age & 0.632 & 0.139 & 4.550 & <.001\\
\bottomrule
\end{tabular}
\end{table}

The fitted regression line is:

\begin{equation}
\label{eq: eq1}
\begin{aligned}
\operatorname{\widehat{chol}} &= 196.97 + 0.63(\operatorname{age})
\end{aligned}
\end{equation}

\underline{\texttt{Interpretation of regression coefficients:}}

The cholesterol level is expected to increase by 0.63 units
for 1-year increase of age.

\hypertarget{answer-to-ques-05}{%
\subsection{Answer To Ques 05}\label{answer-to-ques-05}}

\begin{table}[H]

\caption{\label{tab:Table-07}Estimate of Model Parameters}
\centering
\begin{tabular}[t]{lrrrr}
\toprule
term & estimate & std.error & statistic & p.value\\
\midrule
(Intercept) & 228.352 & 1.097 & 208.102 & <.001\\
dibpatType B & -4.147 & 1.546 & -2.682 & 0.007\\
\bottomrule
\end{tabular}
\end{table}

The fitted regression line is:

\begin{equation}
\label{eq: eq2}
\begin{aligned}
\operatorname{\widehat{chol}} &= 228.35 - 4.15(\operatorname{dibpat}_{\operatorname{Type\ B}})
\end{aligned}
\end{equation}

\underline{\texttt{Interpretation of regression coefficients:}}

\begin{itemize}
\tightlist
\item
  Mean cholesterol level of \texttt{Type\ A\ dibpat} subjects is 228.35.
\item
  Also, mean cholesterol level of \texttt{Type\ A\ dibpat} subjects is 4.15 unit higher than \texttt{Type\ B\ dibpat} subjects.
\end{itemize}

\clearpage

\hypertarget{answer-to-ques-06}{%
\subsection{Answer To Ques 06}\label{answer-to-ques-06}}

\begin{table}[H]

\caption{\label{tab:Table-08}Estimate of Model Parameters}
\centering
\begin{tabular}[t]{lrrrr}
\toprule
term & estimate & std.error & statistic & p.value\\
\midrule
(Intercept) & 235.073 & 2.618 & 89.778 & <.001\\
behpatA2 & -8.150 & 2.883 & -2.827 & 0.005\\
behpatB3 & -10.226 & 2.894 & -3.533 & <.001\\
behpatB4 & -13.122 & 3.491 & -3.759 & <.001\\
\bottomrule
\end{tabular}
\end{table}

The fitted regression line is:

\begin{equation}
\label{eq: eq3}
\begin{aligned}
\operatorname{\widehat{chol}} &= 235.07 - 8.15(\operatorname{behpat}_{\operatorname{A2}}) - 10.23(\operatorname{behpat}_{\operatorname{B3}}) - 13.12(\operatorname{behpat}_{\operatorname{B4}})
\end{aligned}
\end{equation}

\underline{\texttt{Interpretation of regression coefficients:}}

Since \texttt{Type\ A1} Behavior pattern subjects are reference group,

\begin{itemize}
\tightlist
\item
  Mean cholesterol level of \texttt{Type\ A1} subjects is 235.07 unit.
\item
  Mean cholesterol level of \texttt{Type\ A2} subjects is -8.15 unit lower than that of \texttt{Type\ A1} subjects.
\item
  Mean cholesterol level of \texttt{Type\ B3} subjects is -10.23 unit lower than that of \texttt{Type\ A1} subjects.
\item
  Mean cholesterol level of \texttt{Type\ B4} subjects is -13.12 unit lower than that of \texttt{Type\ A1} subjects.
\end{itemize}

\clearpage

\hypertarget{answer-to-ques-07}{%
\subsection{Answer To Ques 07}\label{answer-to-ques-07}}

\begin{table}[H]

\caption{\label{tab:Table-09}Estimate of Model Parameters}
\centering
\begin{tabular}[t]{lrrrr}
\toprule
term & estimate & std.error & statistic & p.value\\
\midrule
(Intercept) & 200.191 & 6.637 & 30.163 & <.001\\
age & 0.600 & 0.140 & 4.302 & <.001\\
dibpatType B & -3.469 & 1.550 & -2.238 & 0.025\\
\bottomrule
\end{tabular}
\end{table}

The fitted regression line is:

\begin{equation}
\label{eq: eq4}
\begin{aligned}
\operatorname{\widehat{chol}} &= 200.19 + 0.6(\operatorname{age}) - 3.47(\operatorname{dibpat}_{\operatorname{Type\ B}})
\end{aligned}
\end{equation}

\underline{\texttt{Interpretation of regression coefficients:}}

Since \texttt{Type\ A1} subjects are reference group,

\begin{itemize}
\tightlist
\item
  Mean cholesterol level of subjects with age 0 and \texttt{Type\ A} pattern is 200.19 unit.
\item
  Mean cholesterol level increases about 0.6 for 1-year increase of age, holding dibpat fixed.
\item
  Mean cholesterol level of \texttt{Type\ B2} subjects is -3.47 unit lower than that of \texttt{Type\ A} dibpat subjects, holding the subject's age fixed.
\end{itemize}

Here, Both the regression coefficients corresponding to age and dibpat (Eq \ref{eq: eq4}) have changed from the case of simple linear regression in Eq \ref{eq: eq1} and Eq \ref{eq: eq2}. And also the value of \(R^2_{adj}\) increased.

\clearpage

\hypertarget{answer-to-ques-08}{%
\subsection{Answer To Ques 08}\label{answer-to-ques-08}}

\begin{table}[H]

\caption{\label{tab:Table-10}Estimate of Model Parameters}
\centering
\begin{tabular}[t]{lrrrr}
\toprule
term & estimate & std.error & statistic & p.value\\
\midrule
(Intercept) & 222.252 & 1.171 & 189.755 & <.001\\
age40 & 0.632 & 0.139 & 4.550 & <.001\\
\bottomrule
\end{tabular}
\end{table}

The new fitted model after subtracting 40 from the variable \texttt{age} is:

\begin{equation}
\label{eq: eq5}
\begin{aligned}
\operatorname{\widehat{chol}} &= 222.25 + 0.63(\operatorname{age40})
\end{aligned}
\end{equation}

\underline{\texttt{Interpretation of regression coefficients:}}

\begin{itemize}
\tightlist
\item
  Since the explanatory variable is age minus 40, we can say, the mean cholesterol level of a 40 year old subject is 222.25.
\item
  Mean cholesterol level is expected to in increase by 0.63 for 1 year increase of age.
\end{itemize}

The main difference of this model (Eq \ref{eq: eq1}) compared to model \ref{eq: eq5} is that, we can interpret the intercept term for this model logically.

\clearpage

\hypertarget{answer-to-ques-09}{%
\subsection{Answer To Ques 09}\label{answer-to-ques-09}}

\begin{table}[H]

\caption{\label{tab:Table-11}Estimate of Model Parameters}
\centering
\begin{tabular}[t]{lrrrr}
\toprule
term & estimate & std.error & statistic & p.value\\
\midrule
(Intercept) & 223.702 & 1.597 & 140.043 & <.001\\
dibpatType B & -3.644 & 2.175 & -1.675 & 0.094\\
smokeYes & 8.728 & 2.189 & 3.988 & <.001\\
dibpatType B:smokeYes & 0.265 & 3.085 & 0.086 & 0.932\\
\bottomrule
\end{tabular}
\end{table}

In this case, the fitted model is:

\begin{equation}
\label{eq: eq6}
\begin{aligned}
\operatorname{\widehat{chol}} &= 223.7 - 3.64(\operatorname{dibpat}_{\operatorname{Type\ B}}) + 8.73(\operatorname{smoke}_{\operatorname{Yes}}) + 0.27(\operatorname{dibpat}_{\operatorname{Type\ B}} \times \operatorname{smoke}_{\operatorname{Yes}})
\end{aligned}
\end{equation}

\underline{\texttt{Interpretation of regression coefficients:}}

Here the reference groups are \texttt{Non-smoker} and dibpat \texttt{Type\ A} subjects. So,

\begin{itemize}
\tightlist
\item
  The mean cholesterol level of Type A non smoker subject is 223.7 unit.
\item
  Among the non-smokers, mean cholesterol level of dibpat \texttt{Type\ B} subjects is 3.64 unit lower compared to \texttt{Type\ A} subject.
\item
  Among the \texttt{Type\ A} subjects, mean cholesterol level of smokers is 8.73 unit higher than that of non-smokers.
\item
  Difference of mean cholesterol level between smokers and non-smokers is 0.27 unit higher in \texttt{Type\ B} dibpat subjects compared to that of \texttt{Type\ A} dibpat subjects.
\end{itemize}

\clearpage

\hypertarget{r-code}{%
\subsection{R-code}\label{r-code}}

\begin{Shaded}
\begin{Highlighting}[]
\NormalTok{knitr}\SpecialCharTok{::}\NormalTok{opts\_chunk}\SpecialCharTok{$}\FunctionTok{set}\NormalTok{(}
  \AttributeTok{echo =} \ConstantTok{FALSE}\NormalTok{,}
  \AttributeTok{message =} \ConstantTok{FALSE}\NormalTok{,}
  \AttributeTok{warning =} \ConstantTok{FALSE}
\NormalTok{)}

\DocumentationTok{\#\# {-}{-}{-}{-}{-}{-}{-} package setup {-}{-}{-}{-}{-}{-}{-}{-}{-}{-}{-}{-}{-}{-}{-}{-}{-}{-}{-}{-}{-}{-}{-}{-}}

\FunctionTok{library}\NormalTok{(dplyr)}
\FunctionTok{library}\NormalTok{(purrr)}
\FunctionTok{library}\NormalTok{(knitr)}
\FunctionTok{library}\NormalTok{(broom)}
\FunctionTok{library}\NormalTok{(tableone)}
\FunctionTok{library}\NormalTok{(kableExtra)}
\FunctionTok{library}\NormalTok{(equatiomatic)}


\DocumentationTok{\#\# {-}{-}{-}{-}{-}{-}{-} data setup {-}{-}{-}{-}{-}{-}{-}{-}{-}{-}{-}{-}{-}{-}{-}{-}{-}{-}{-}{-}{-}{-}{-}{-}}

\FunctionTok{load}\NormalTok{(here}\SpecialCharTok{::}\FunctionTok{here}\NormalTok{(}\StringTok{"data"}\NormalTok{, }\StringTok{"wcgs.Rdata"}\NormalTok{))}
\NormalTok{sid }\OtherTok{\textless{}{-}} \DecValTok{011}
\FunctionTok{set.seed}\NormalTok{(sid)}
\NormalTok{mydat }\OtherTok{\textless{}{-}} \FunctionTok{sample\_n}\NormalTok{(wcgs, }\AttributeTok{size =} \FunctionTok{n}\NormalTok{(), }\AttributeTok{replace =} \ConstantTok{TRUE}\NormalTok{)}


\DocumentationTok{\#\# {-}{-}{-}{-}{-}{-}{-} utility functions {-}{-}{-}{-}{-}{-}{-}{-}{-}{-}{-}{-}{-}{-}{-}{-}{-}{-}{-}{-}{-}{-}{-}{-}}

\NormalTok{kab\_tab }\OtherTok{\textless{}{-}} \ControlFlowTok{function}\NormalTok{(tab, ...) \{}
\NormalTok{  knitr}\SpecialCharTok{::}\FunctionTok{kable}\NormalTok{(tab,}
               \AttributeTok{format =} \StringTok{"latex"}\NormalTok{,}
               \AttributeTok{booktabs =} \ConstantTok{TRUE}\NormalTok{,}
               \AttributeTok{digits =} \DecValTok{3}\NormalTok{,}
\NormalTok{               ...)}
\NormalTok{\}}

\NormalTok{p\_format }\OtherTok{\textless{}{-}} \ControlFlowTok{function}\NormalTok{(pval) \{}
  \FunctionTok{ifelse}\NormalTok{(pval }\SpecialCharTok{\textless{}}\NormalTok{ .}\DecValTok{001}\NormalTok{, }\StringTok{"\textless{}.001"}\NormalTok{, }\FunctionTok{as.character}\NormalTok{(}\FunctionTok{round}\NormalTok{(pval, }\DecValTok{3}\NormalTok{)))}
\NormalTok{\}}

\NormalTok{mod\_tab }\OtherTok{\textless{}{-}} \ControlFlowTok{function}\NormalTok{(mod, ...) \{}
\NormalTok{  mod }\SpecialCharTok{\%\textgreater{}\%}
    \FunctionTok{tidy}\NormalTok{() }\SpecialCharTok{\%\textgreater{}\%}
    \FunctionTok{mutate}\NormalTok{(}\AttributeTok{p.value =} \FunctionTok{p\_format}\NormalTok{(p.value)) }\SpecialCharTok{\%\textgreater{}\%}
    \FunctionTok{kab\_tab}\NormalTok{(}\AttributeTok{align =} \StringTok{"lrrrr"}\NormalTok{,}
            \AttributeTok{caption =} \StringTok{"Estimate of Model Parameters"}\NormalTok{, ...) }\SpecialCharTok{\%\textgreater{}\%}
    \FunctionTok{kable\_styling}\NormalTok{(}\AttributeTok{latex\_options =} \StringTok{"HOLD\_position"}\NormalTok{)}
\NormalTok{\}}

\NormalTok{reg\_eq }\OtherTok{\textless{}{-}} \ControlFlowTok{function}\NormalTok{(mod, ref, ...) \{}
  \FunctionTok{extract\_eq}\NormalTok{(mod,}
             \AttributeTok{use\_coefs =} \ConstantTok{TRUE}\NormalTok{,}
             \AttributeTok{intercept =} \StringTok{"beta"}\NormalTok{,}
             \AttributeTok{wrap =} \ConstantTok{TRUE}\NormalTok{,}
             \AttributeTok{label =} \FunctionTok{paste0}\NormalTok{(}\StringTok{"eq"}\NormalTok{,ref),}
\NormalTok{             ...)}
\NormalTok{\}}

\NormalTok{params }\OtherTok{\textless{}{-}} \ControlFlowTok{function}\NormalTok{(mod, param, }\AttributeTok{dec =} \DecValTok{2}\NormalTok{) \{}
  \FunctionTok{round}\NormalTok{(mod}\SpecialCharTok{$}\NormalTok{coefficients[[param]], dec)}
\NormalTok{\}}











\DocumentationTok{\#\# {-}{-}{-}{-}{-}{-}{-} Code for Question{-}01 {-}{-}{-}{-}{-}{-}{-}{-}{-}{-}{-}{-}{-}{-}{-}{-}{-}{-}{-}{-}{-}{-}{-}{-}}

\NormalTok{tab }\OtherTok{\textless{}{-}} \FunctionTok{CreateTableOne}\NormalTok{(}
  \AttributeTok{data =}\NormalTok{ mydat,}
  \AttributeTok{vars =} \FunctionTok{c}\NormalTok{(}\StringTok{"age"}\NormalTok{, }\StringTok{"bmi"}\NormalTok{, }\StringTok{"arcus"}\NormalTok{, }\StringTok{"behpat"}\NormalTok{, }\StringTok{"chd69"}\NormalTok{),}
  \AttributeTok{factorVars =} \StringTok{"arcus"}\NormalTok{,}
  \AttributeTok{addOverall =} \ConstantTok{TRUE}
\NormalTok{)}

\NormalTok{tab\_p }\OtherTok{\textless{}{-}} \FunctionTok{print}\NormalTok{(tab, }\AttributeTok{showAllLevels =} \ConstantTok{TRUE}\NormalTok{, }\AttributeTok{printToggle =} \ConstantTok{FALSE}\NormalTok{)}

\FunctionTok{kab\_tab}\NormalTok{(tab\_p,}
        \AttributeTok{caption =} \StringTok{"Descriptive Statistics"}\NormalTok{) }\SpecialCharTok{\%\textgreater{}\%}
  \FunctionTok{kable\_styling}\NormalTok{(}\AttributeTok{latex\_options =} \StringTok{"hold\_position"}\NormalTok{)}


\DocumentationTok{\#\# {-}{-}{-}{-}{-}{-}{-} Code for Question{-}02 {-}{-}{-}{-}{-}{-}{-}{-}{-}{-}{-}{-}{-}{-}{-}{-}{-}{-}{-}{-}{-}{-}{-}{-}}

\NormalTok{arcus }\OtherTok{\textless{}{-}} \FunctionTok{t.test}\NormalTok{(chol }\SpecialCharTok{\textasciitilde{}}\NormalTok{ arcus, }\AttributeTok{data =}\NormalTok{ mydat) }\SpecialCharTok{\%\textgreater{}\%}
  \FunctionTok{tidy}\NormalTok{() }\SpecialCharTok{\%\textgreater{}\%}
  \FunctionTok{mutate}\NormalTok{(}\AttributeTok{term =} \StringTok{"arcus"}\NormalTok{) }\SpecialCharTok{\%\textgreater{}\%}
  \FunctionTok{select}\NormalTok{(term, estimate, statistic, p.value)}

\NormalTok{mydat }\SpecialCharTok{\%\textgreater{}\%}
  \FunctionTok{select}\NormalTok{(age, bmi) }\SpecialCharTok{\%\textgreater{}\%}
  \FunctionTok{map}\NormalTok{(}\SpecialCharTok{\textasciitilde{}} \FunctionTok{cor.test}\NormalTok{(}\AttributeTok{x =}\NormalTok{ .x, }\AttributeTok{y =}\NormalTok{ mydat}\SpecialCharTok{$}\NormalTok{chol)) }\SpecialCharTok{\%\textgreater{}\%}
  \FunctionTok{map\_dfr}\NormalTok{(broom}\SpecialCharTok{::}\NormalTok{tidy, }\AttributeTok{.id =} \StringTok{"term"}\NormalTok{) }\SpecialCharTok{\%\textgreater{}\%}
  \FunctionTok{select}\NormalTok{(term}\SpecialCharTok{:}\NormalTok{p.value) }\SpecialCharTok{\%\textgreater{}\%}
  \FunctionTok{bind\_rows}\NormalTok{(arcus) }\SpecialCharTok{\%\textgreater{}\%}
  \FunctionTok{mutate}\NormalTok{(}\AttributeTok{p.value =} \FunctionTok{p\_format}\NormalTok{(p.value)) }\SpecialCharTok{\%\textgreater{}\%}
  \FunctionTok{kab\_tab}\NormalTok{(}\AttributeTok{align =} \StringTok{"lrrr"}\NormalTok{,}
          \AttributeTok{caption =} \StringTok{"Correlation Test (age, bmi) and T{-}test (arcus) with chol"}\NormalTok{) }\SpecialCharTok{\%\textgreater{}\%}
  \FunctionTok{kable\_styling}\NormalTok{(}\AttributeTok{latex\_options =} \StringTok{"hold\_position"}\NormalTok{)}


\FunctionTok{anova}\NormalTok{(}\FunctionTok{lm}\NormalTok{(chol }\SpecialCharTok{\textasciitilde{}}\NormalTok{ behpat, }\AttributeTok{data =}\NormalTok{ mydat)) }\SpecialCharTok{\%\textgreater{}\%}
  \FunctionTok{tidy}\NormalTok{() }\SpecialCharTok{\%\textgreater{}\%}
  \FunctionTok{mutate}\NormalTok{(}\AttributeTok{p.value =} \FunctionTok{p\_format}\NormalTok{(p.value)) }\SpecialCharTok{\%\textgreater{}\%}
  \FunctionTok{kab\_tab}\NormalTok{(}\AttributeTok{align =} \StringTok{"lrrrrr"}\NormalTok{,}
          \AttributeTok{caption =} \StringTok{"One Way Analysis of Variance for chol on behpat"}\NormalTok{) }\SpecialCharTok{\%\textgreater{}\%}
  \FunctionTok{kable\_styling}\NormalTok{(}\AttributeTok{latex\_options =} \StringTok{"HOLD\_position"}\NormalTok{)}


\DocumentationTok{\#\# {-}{-}{-}{-}{-}{-}{-} Code for Question{-}03 {-}{-}{-}{-}{-}{-}{-}{-}{-}{-}{-}{-}{-}{-}{-}{-}{-}{-}{-}{-}{-}{-}{-}{-}}

\NormalTok{df\_cont }\OtherTok{\textless{}{-}}\NormalTok{ mydat }\SpecialCharTok{\%\textgreater{}\%}\NormalTok{ janitor}\SpecialCharTok{::}\FunctionTok{tabyl}\NormalTok{(behpat, chd69)}

\NormalTok{tab\_cont }\OtherTok{\textless{}{-}}\NormalTok{ df\_cont }\SpecialCharTok{\%\textgreater{}\%}
  \FunctionTok{kab\_tab}\NormalTok{() }\SpecialCharTok{\%\textgreater{}\%}
  \FunctionTok{add\_header\_above}\NormalTok{(}\AttributeTok{header =} \FunctionTok{c}\NormalTok{(}\StringTok{" "} \OtherTok{=} \DecValTok{1}\NormalTok{, }\StringTok{"CHD"} \OtherTok{=} \DecValTok{2}\NormalTok{))}

\NormalTok{tab\_chi }\OtherTok{\textless{}{-}}\NormalTok{ df\_cont }\SpecialCharTok{\%\textgreater{}\%}
\NormalTok{  janitor}\SpecialCharTok{::}\FunctionTok{chisq.test}\NormalTok{() }\SpecialCharTok{\%\textgreater{}\%}
  \FunctionTok{tidy}\NormalTok{() }\SpecialCharTok{\%\textgreater{}\%}
  \FunctionTok{select}\NormalTok{(statistic, p.value) }\SpecialCharTok{\%\textgreater{}\%}
  \FunctionTok{mutate}\NormalTok{(}\AttributeTok{p.value =} \FunctionTok{p\_format}\NormalTok{(p.value)) }\SpecialCharTok{\%\textgreater{}\%}
  \FunctionTok{kab\_tab}\NormalTok{()}

\NormalTok{tab\_side }\OtherTok{\textless{}{-}} \FunctionTok{c}\NormalTok{(}
    \StringTok{"}\SpecialCharTok{\textbackslash{}\textbackslash{}}\StringTok{begin\{table\}[H]}
\StringTok{      }\SpecialCharTok{\textbackslash{}\textbackslash{}}\StringTok{begin\{minipage\}\{.5}\SpecialCharTok{\textbackslash{}\textbackslash{}}\StringTok{linewidth\}}
\StringTok{      }\SpecialCharTok{\textbackslash{}\textbackslash{}}\StringTok{centering"}\NormalTok{,}
\NormalTok{      tab\_cont,}
      \StringTok{"}\SpecialCharTok{\textbackslash{}\textbackslash{}}\StringTok{caption\{Contingency table\}}
\StringTok{      }\SpecialCharTok{\textbackslash{}\textbackslash{}}\StringTok{label\{Table{-}04\}}
\StringTok{    }\SpecialCharTok{\textbackslash{}\textbackslash{}}\StringTok{end\{minipage\}\%}
\StringTok{      }\SpecialCharTok{\textbackslash{}\textbackslash{}}\StringTok{begin\{minipage\}\{.5}\SpecialCharTok{\textbackslash{}\textbackslash{}}\StringTok{linewidth\}}
\StringTok{        }\SpecialCharTok{\textbackslash{}\textbackslash{}}\StringTok{centering"}\NormalTok{,}
\NormalTok{    tab\_chi,}
    \StringTok{"}\SpecialCharTok{\textbackslash{}\textbackslash{}}\StringTok{caption\{chi{-}square test\}}
\StringTok{    }\SpecialCharTok{\textbackslash{}\textbackslash{}}\StringTok{label\{Table{-}05\}}
\StringTok{    }\SpecialCharTok{\textbackslash{}\textbackslash{}}\StringTok{end\{minipage\}}
\StringTok{      }\SpecialCharTok{\textbackslash{}\textbackslash{}}\StringTok{end\{table\}"}
\NormalTok{  )}


\DocumentationTok{\#\# {-}{-}{-}{-}{-}{-}{-} Code for Question{-}04 {-}{-}{-}{-}{-}{-}{-}{-}{-}{-}{-}{-}{-}{-}{-}{-}{-}{-}{-}{-}{-}{-}{-}{-}}

\NormalTok{m1 }\OtherTok{\textless{}{-}} \FunctionTok{lm}\NormalTok{(chol }\SpecialCharTok{\textasciitilde{}}\NormalTok{ age, }\AttributeTok{data =}\NormalTok{ mydat)}
\NormalTok{m1 }\SpecialCharTok{\%\textgreater{}\%} \FunctionTok{mod\_tab}\NormalTok{()}

\FunctionTok{reg\_eq}\NormalTok{(m1, }\DecValTok{1}\NormalTok{)}

\DocumentationTok{\#\# {-}{-}{-}{-}{-}{-}{-} Code for Question{-}05 {-}{-}{-}{-}{-}{-}{-}{-}{-}{-}{-}{-}{-}{-}{-}{-}{-}{-}{-}{-}{-}{-}{-}{-}}

\NormalTok{m2 }\OtherTok{\textless{}{-}} \FunctionTok{lm}\NormalTok{(chol }\SpecialCharTok{\textasciitilde{}}\NormalTok{ dibpat, }\AttributeTok{data =}\NormalTok{ mydat)}
\NormalTok{m2 }\SpecialCharTok{\%\textgreater{}\%} \FunctionTok{mod\_tab}\NormalTok{()}

\FunctionTok{reg\_eq}\NormalTok{(m2, }\DecValTok{2}\NormalTok{)}

\DocumentationTok{\#\# {-}{-}{-}{-}{-}{-}{-} Code for Question{-}05 {-}{-}{-}{-}{-}{-}{-}{-}{-}{-}{-}{-}{-}{-}{-}{-}{-}{-}{-}{-}{-}{-}{-}{-}}

\NormalTok{m3 }\OtherTok{\textless{}{-}} \FunctionTok{lm}\NormalTok{(chol }\SpecialCharTok{\textasciitilde{}}\NormalTok{ behpat, }\AttributeTok{data =}\NormalTok{ mydat)}
\NormalTok{m3 }\SpecialCharTok{\%\textgreater{}\%} \FunctionTok{mod\_tab}\NormalTok{()}

\FunctionTok{reg\_eq}\NormalTok{(m3, }\DecValTok{3}\NormalTok{)}

\DocumentationTok{\#\# {-}{-}{-}{-}{-}{-}{-} Code for Question{-}08 {-}{-}{-}{-}{-}{-}{-}{-}{-}{-}{-}{-}{-}{-}{-}{-}{-}{-}{-}{-}{-}{-}{-}{-}}

\NormalTok{m4 }\OtherTok{\textless{}{-}} \FunctionTok{lm}\NormalTok{(chol }\SpecialCharTok{\textasciitilde{}}\NormalTok{ age }\SpecialCharTok{+}\NormalTok{ dibpat, }\AttributeTok{data =}\NormalTok{ mydat)}
\NormalTok{m4 }\SpecialCharTok{\%\textgreater{}\%} \FunctionTok{mod\_tab}\NormalTok{()}

\FunctionTok{reg\_eq}\NormalTok{(m4, }\DecValTok{4}\NormalTok{)}

\DocumentationTok{\#\# {-}{-}{-}{-}{-}{-}{-} Code for Question{-}08 {-}{-}{-}{-}{-}{-}{-}{-}{-}{-}{-}{-}{-}{-}{-}{-}{-}{-}{-}{-}{-}{-}{-}{-}}

\NormalTok{mydat }\SpecialCharTok{\%\textgreater{}\%}
  \FunctionTok{mutate}\NormalTok{(}\AttributeTok{age40 =}\NormalTok{ age }\SpecialCharTok{{-}} \DecValTok{40}\NormalTok{) }\SpecialCharTok{\%\textgreater{}\%}
  \FunctionTok{lm}\NormalTok{(chol }\SpecialCharTok{\textasciitilde{}}\NormalTok{ age40, }\AttributeTok{data =}\NormalTok{ .) }\OtherTok{{-}\textgreater{}}\NormalTok{ m5}

\NormalTok{m5 }\SpecialCharTok{\%\textgreater{}\%} \FunctionTok{mod\_tab}\NormalTok{()}

\FunctionTok{reg\_eq}\NormalTok{(m5, }\DecValTok{5}\NormalTok{)}
\DocumentationTok{\#\# {-}{-}{-}{-}{-}{-}{-} Code for Question{-}08 {-}{-}{-}{-}{-}{-}{-}{-}{-}{-}{-}{-}{-}{-}{-}{-}{-}{-}{-}{-}{-}{-}{-}{-}}

\NormalTok{m6 }\OtherTok{\textless{}{-}} \FunctionTok{lm}\NormalTok{(chol }\SpecialCharTok{\textasciitilde{}}\NormalTok{ dibpat }\SpecialCharTok{*}\NormalTok{ smoke, }\AttributeTok{data =}\NormalTok{ mydat)}
\NormalTok{m6 }\SpecialCharTok{\%\textgreater{}\%} \FunctionTok{mod\_tab}\NormalTok{()}
\FunctionTok{reg\_eq}\NormalTok{(m6, }\DecValTok{6}\NormalTok{)}
\end{Highlighting}
\end{Shaded}


\end{document}
